% ID: SEL-D0084-OA-E02
% Tanda S084.
\textbf{a)} La derivada $f'(x)=x+2$ es negativa para $x<-2$ y positiva para $x>-2$. Por tanto, $f$ decrece en $(-\infty,-2)$ y crece en $(-2,\infty)$, con un mínimo relativo en $x=-2$. Como $g'(x)=2>0$ para todo $x$, $g$ es estrictamente creciente en $\mathbb R$.

\textbf{b)} Solo $f$ tiene derivada nula, exactamente en $x=-2$. La derivada de $g$ vale 2 en todos los puntos.

\textbf{c)} Una función con derivada constante 2 tiene la forma $g(x)=2x+C$, con $C$ constante. Es un polinomio de primer grado, pues el coeficiente de $x$ es distinto de cero. La función $f$ tiene derivada no constante y corresponde a una cuadrática, $f(x)=x^2/2+2x+K$.
